% --- INICIO DE PLANTILLA PERSONALIZADA ---
\documentclass[oneside,11pt]{Latex/Classes/PhDthesisPSnPDF}

% --- Paquetes Originales ---
\usepackage{blindtext}
\usepackage{actuarialangle} 
\usepackage{tabularx}
\usepackage{float}
\usepackage{multirow, array}
\usepackage[usenames,dvipsnames,svgnames,table]{xcolor}
\usepackage{longtable}
\usepackage{latexsym,amsmath,amssymb,amsfonts}
\usepackage{enumitem}
\usepackage{xparse}
\usepackage{booktabs}
\usepackage[round, sort, numbers]{natbib}  
\usepackage{adjustbox}
\usepackage[utf8]{inputenc}
\usepackage{caption}
\usepackage{graphicx}
\usepackage{hyperref}
\usepackage{cleveref}

% Definición de colores
\definecolor{miblue}{RGB}{68, 114, 196}

% --- CAMBIO CRÍTICO 1: Usar \input para macros en el preámbulo ---
% \include da problemas antes del begin document. \input es seguro.
\input{Latex/Macros/MacroFile1}

% --- Definiciones de seguridad (por si la clase falla al cargar alguna variable) ---
\providecommand{\facultad}[1]{\def\lafacultad{#1}}
\providecommand{\escudofacultad}[1]{\def\elescudo{#1}}
\providecommand{\degree}[1]{\def\elgrado{#1}}
\providecommand{\director}[1]{\def\eldirector{#1}}
\providecommand{\lugar}[1]{\def\ellugar{#1}}
\providecommand{\degreedate}[1]{\def\lafecha{#1}}

% --- Configuración para bloques de código de R (Pandoc) ---

%%%%%%%%%%%%%%%%%%%%%%%%%%%%%%%%%%%%%%%%%%%%%%%%%%%%%%%%%%%%%%%%%%%%%%%%%%%%%%%%
%                         DATOS (Conectados a RMarkdown)                       %
%%%%%%%%%%%%%%%%%%%%%%%%%%%%%%%%%%%%%%%%%%%%%%%%%%%%%%%%%%%%%%%%%%%%%%%%%%%%%%%%
\title{De la estadistica a la victoria}
\author{David Linarez}

% Lógica para Facultad: Si la defines en el YAML la usa, si no, usa la default

\facultad{Facultad de Ciencias Económicas y Sociales \\ Escuela de
Estadistica y Ciencias Actuariales}
  \escudofacultad{Latex/Classes/Escudos/faces}

\degree{Computación I}
\director{Jesus Ochoa \\ Oliver Triveño}
\degreedate{Abril 2026} 
\lugar{Caracas}

\portadatrue 

% Metadatos PDF
\keywords{}
\subject{}

% --- CORRECCIÓN PARA PANDOC NUEVO (Imágenes) ---
\makeatletter
\providecommand{\pandocbounded}[1]{#1}
\makeatother


%%%%%%%%%%%%%%%%%%%%%%%%%%%%%%%%%%%%%%%%%%%%%%%%%%%%%
%                   DOCUMENTO                       %
%%%%%%%%%%%%%%%%%%%%%%%%%%%%%%%%%%%%%%%%%%%%%%%%%%%%%
\begin{document}

\maketitle

%%%%%%%%%%%%%%%%%%%%%%%%%%%%%%%%%%%%%%%%%%%%%%%%%%%%%
%                  PRÓLOGO                          %
%%%%%%%%%%%%%%%%%%%%%%%%%%%%%%%%%%%%%%%%%%%%%%%%%%%%%
\frontmatter

% Puedes agregar dedicatorias desde el YAML si quieres, o dejarlas fijas aquí

%%%%%%%%%%%%%%%%%%%%%%%%%%%%%%%%%%%%%%%%%%%%%%%%%%%%%
%                   ÍNDICES                         %
%%%%%%%%%%%%%%%%%%%%%%%%%%%%%%%%%%%%%%%%%%%%%%%%%%%%%
\setcounter{secnumdepth}{3}
\setcounter{tocdepth}{3}

\tableofcontents



\cleardoublepage

  \cleardoublepage       % Empieza en página derecha (estándar en tesis)
  \chapter*{Resumen}     % Crea el título "Resumen" sin numeración (Capítulo *)
  \addcontentsline{toc}{chapter}{Resumen} % (Opcional) Lo añade al índice
  
  El presente trabajo tiene como objetivo determinar la formación ideal
de un equipo a través de modelos estadísticos. Mediante el procesamiento
de las estadísticas base de cada categoría, se busca desarrollar una
estructura equilibrada y competitiva que maximice el rendimiento
colectivo.             % Aquí se pega el texto que escribiste en el YAML
  
  \cleardoublepage

%%%%%%%%%%%%%%%%%%%%%%%%%%%%%%%%%%%%%%%%%%%%%%%%%%%%%
%                   CONTENIDO (El Corazón)          %
%%%%%%%%%%%%%%%%%%%%%%%%%%%%%%%%%%%%%%%%%%%%%%%%%%%%%
\mainmatter
\def\baselinestretch{1.5}

% --- CAMBIO CRÍTICO 2: Aquí RMarkdown inyecta todo tu texto ---
\chapter*{Introducción}

Históricamente, la experiencia de muchos jugadores de Pokémon comenzó
con una estrategia novata, que consiste en fortalecer a un único
pokemon, casi siempre el pokemon ``favorito'' hasta que su nivel
superara cualquier obstáculo por fuerza bruta. Sin embargo, al analizar
el juego desde una perspectiva técnica y competitiva, surge una
interrogante fundamental ¿cómo se construye una formación verdaderamente
equilibrada? El presente trabajo trasciende la preferencia individual
para proponer un modelo estadístico de selección. El objetivo es
ofrecer, tanto a jugadores veteranos como a novatos, una metodología
clara para estructurar equipos basados en sus tipologías predilectas,
garantizando que la sinergia y las estadísticas base trabajen en favor
de un rendimiento óptimo en combate.

\chapter{Planteamiento del Problema}

Desde su creación, la franquicia Pokémon ha evolucionado hacia un
ecosistema competitivo de alta complejidad donde la victoria no depende
del azar, sino de la optimización matemática de recursos. Para un
entrenador o analista, el desafío principal radica en la selección de un
equipo equilibrado, enfrentándose a una variabilidad extrema en las
combinaciones de tipos y naturalezas que definen el rendimiento en
combate.

Desde una perspectiva de análisis estadístico, el problema no consiste
solo en elegir criaturas populares, sino en identificar los puntos altos
de las estadísticas base (HP, Atk, Def, Sp. Atk, Sp. Def y Speed) según
la tipología elemental. Factores como la especialización de roles
(Atacantes, Tanques, Muros) introducen variables que dificultan la
creación de un estándar de ``equipo ideal'' sin un estudio previo de los
datos.

La ausencia de un análisis descriptivo que cuantifique cómo los tipos
elementales condicionan la esperanza de desempeño en cada estadística
genera ineficiencias en la toma de decisiones tácticas y en la
construcción de estrategias competitivas a largo plazo.

Dado que un buen equipo depende de la estabilidad y especializaci´on de
sus pokemones, se plantea las siguientes preguntas de investigaci´on::

\textbf{1.} ¿Qué tipos ofrecen las estadísticas más ofensivas?

\textbf{2.} ¿Qué tipos ofrecen las estadísticas más defensivas?

\textbf{3.} ¿Qué tipos ofrecen las estadísticas más veloces?

\textbf{4.} ¿Cuales son los tipos de roles basados en las estadisticas
de cada tipo pokemon?

\textbf{5.} ¿Cuales son las debilidades de cada tipo pokemon?

\section{Justificación}

La presente investigación se justifica en la necesidad de
profesionalizar los criterios de selección en el ámbito competitivo de
Pokémon. Ante un universo de más de mil especies, la toma de decisiones
suele basarse en la intuición, lo que genera desequilibrios en el
combate.

Desde un punto de vista estadístico, este trabajo es relevante porque
permite aplicar medidas de dispersión para reducir la incertidumbre del
jugador. Al identificar qué tipos son más consistentes y cuáles más
erráticos, se ofrece una guía técnica que beneficia tanto a jugadores
principiantes, facilitando su curva de aprendizaje, como a competidores
avanzados que buscan optimizar sus estrategias mediante el análisis de
datos.\\

\section{Objetivos}

\subsection{Objetivo General}
\begin{itemize}
  \item{Analizar estadísticamente la distribución de las estadísticas base (HP, Atk, Def, Sp. Atk, Sp. Def y Speed) de cada tipologia pokemon para determinar patrones de especialización por rol, con el fin de fundamentar la toma de decisiones estratégicas en la construcción de equipos competitivos equilibrados.}
\end{itemize}

\subsection{Objetivos Específicos}
\begin{itemize}
  \item{Identificar las medidas de tendencia central (media) de cada estadística base (HP, Atk, Def, Sp. Atk, Sp. Def y Speed) desglosadas por tipo de Pokémon para establecer un perfil promedio por categoría.}
  \item{Calcular las medidas de dispersión (varianza y desviación estándar) especiales por cada tipo, con el fin de determinar qué categorías son más consistentes o especialistas en sus roles de combate.}
  \item{Categorizar a la tipologia pokemon en roles funcionales (Atacantes, Tanques o Muros).}
  \item{Evaluar el impacto de las debilidades y resistencias elementales en la estabilidad estructural de un equipo}
\end{itemize}

\begin{table}[ht]
\centering
\caption{Variables consideradas en el estudio} 
\resizebox{12cm}{!} {
\begin{tabular}{cll}
  \hline
  \textbf{TIPO} & \textbf{VARIABLE} & \textbf{DESCRIPCIÓN} \\ 
  \hline
  Categórica & Type\_1 & Tipo elemental principal \\ 
  Numérica & HP & Valor numerico que representa la capacidad de vida del pokemon. \\ 
  Numérica & Attack  & Valor numerico que representa la capacidad de ataque ofensiva del pokemon. \\ 
  Numérica & Defense & Valor numerico que representa la capacidad de defensa ofensiva del pokemon \\ 
  Numérica & Sp.Atk & Valor numerico que representa la capacidad de ataque especial del pokemon. \\ 
  Numérica & Sp.Def & Valor numerico que representa la capacidad de defensa especial del pokemon.\\ 
  Numérica & Speed & Valor numerico que representa la velocidad del pokemon. \\ 
  
\end{tabular}
}
\end{table}

\section{Variables}

Las variables principales en estudio (Variables Dependientes/Objetivo)
son las estadísticas base (HP, Attack, Defense, Sp. Atk, Sp. Def y
Speed), las cuales representan el potencial competitivo de cada tipo de
pokemon. Como variable explicativa o independiente, se encuentra la
variable categorica Type\_1.

\chapter{Marco Teórico}

\section{Bases Teóricas}

La franquicia pokemon ha trascendido el ambito del entretenimiento para
convertirse en un objeto de estudio. El nucleo funcional de este
ecosistema reside en las estadisticas base (base stats), las cuales
constituyen las variables intrínsecas que determinan el potencial de
desempeño de cada pokemon.

A diferencia de otras variables modificables (como el entrenamiento o
los objetos), las estadísticas base son constantes universales para cada
especie. Esto permite tratarlas como datos faciles de ver, ideales para
un analisis estadistico que busque identificar patrones de
especialización. El equilibrio de un equipo competitivo no depende
unicamente de la potencia individual, sino de la estratificación de
roles basada en la distribución de estos valores numéricos.\\

\subsection{Categorización de Roles Funcionales}

Desde una perspectiva analítica, los roles se definen mediante la
distribución de los cuantiles de las estadísticas base:

\begin{itemize}
 \item {Atacantes: Pokemones con altos valores situados en las variables de Ataque, Ataque Especial y Velocidad, sacrificando generalmente sus capacidades defensivas.}
 \item {Tanques: Pokemones que presentan un equilibrio estadístico, manteniendo una media elevada tanto en puntos de salud (HP) como en estadísticas ofensivas, permitiendo el intercambio de daño prolongado.}
 \item {Muros: Categoría definida por una baja dispersión en valores altos de Defensa y Defensa Especial, cuya función principal es la absorción de daño y la interrupción del ritmo del oponente.} 
 \end{itemize}

\subsection{Estabilidad Estructural y Análisis de Vulnerabilidad}

Se define como la capacidad de un equipo para mitigar la probabilidad de
daño efectivo mediante la diversificación de tipos. En términos
estadísticos, un equipo con alta estabilidad estructural es aquel que
presenta un balance neto positivo entre debilidades y resistencias,
reduciendo la varianza del daño recibido ante diferentes configuraciones
de rivales.\\

\begin{table}[!h]
\centering
\caption{\label{tab:unnamed-chunk-1}Matriz de Efectividades: Acero a Hielo (Defensores)}
\centering
\resizebox{\ifdim\width>\linewidth\linewidth\else\width\fi}{!}{
\fontsize{10}{12}\selectfont
\begin{tabular}[t]{>{}lrrrrrrrrr}
\toprule
\textbf{ } & \textbf{Acero} & \textbf{Agua} & \textbf{Bicho} & \textbf{Dragón} & \textbf{Eléctrico} & \textbf{Fantasma} & \textbf{Fuego} & \textbf{Hada} & \textbf{Hielo}\\
\midrule
\textbf{Acero} & 0.5 & 0.5 & 1.0 & 1.0 & 0.5 & 1.0 & 0.5 & 2.0 & 2.0\\
\textbf{Agua} & 1.0 & 0.5 & 1.0 & 0.5 & 1.0 & 1.0 & 2.0 & 1.0 & 1.0\\
\textbf{Bicho} & 0.5 & 1.0 & 1.0 & 1.0 & 1.0 & 0.5 & 0.5 & 0.5 & 1.0\\
\textbf{Dragón} & 0.5 & 1.0 & 1.0 & 2.0 & 1.0 & 1.0 & 1.0 & 0.0 & 1.0\\
\textbf{Eléctrico} & 1.0 & 2.0 & 1.0 & 0.5 & 0.5 & 1.0 & 1.0 & 1.0 & 1.0\\
\addlinespace
\textbf{Fantasma} & 1.0 & 1.0 & 1.0 & 1.0 & 1.0 & 2.0 & 1.0 & 1.0 & 1.0\\
\textbf{Fuego} & 2.0 & 0.5 & 2.0 & 0.5 & 1.0 & 1.0 & 0.5 & 1.0 & 2.0\\
\textbf{Hada} & 0.5 & 1.0 & 1.0 & 2.0 & 1.0 & 1.0 & 0.5 & 1.0 & 1.0\\
\textbf{Hielo} & 0.5 & 0.5 & 1.0 & 2.0 & 1.0 & 1.0 & 0.5 & 1.0 & 0.5\\
\textbf{Lucha} & 2.0 & 1.0 & 0.5 & 1.0 & 1.0 & 0.0 & 1.0 & 0.5 & 2.0\\
\addlinespace
\textbf{Normal} & 0.5 & 1.0 & 1.0 & 1.0 & 1.0 & 0.0 & 1.0 & 1.0 & 1.0\\
\textbf{Planta} & 0.5 & 2.0 & 0.5 & 0.5 & 1.0 & 1.0 & 0.5 & 1.0 & 1.0\\
\textbf{Psíquico} & 0.5 & 1.0 & 1.0 & 1.0 & 1.0 & 1.0 & 1.0 & 1.0 & 1.0\\
\textbf{Roca} & 0.5 & 1.0 & 2.0 & 1.0 & 1.0 & 1.0 & 2.0 & 1.0 & 2.0\\
\textbf{Siniestro} & 1.0 & 1.0 & 1.0 & 1.0 & 1.0 & 2.0 & 1.0 & 0.5 & 1.0\\
\addlinespace
\textbf{Tierra} & 2.0 & 1.0 & 0.5 & 1.0 & 2.0 & 1.0 & 2.0 & 1.0 & 1.0\\
\textbf{Veneno} & 0.0 & 1.0 & 1.0 & 1.0 & 1.0 & 0.5 & 1.0 & 2.0 & 1.0\\
\textbf{Volador} & 0.5 & 1.0 & 2.0 & 1.0 & 0.5 & 1.0 & 1.0 & 1.0 & 1.0\\
\bottomrule
\end{tabular}}
\end{table}

\begin{table}[!h]
\centering
\caption{\label{tab:unnamed-chunk-1}Matriz de Efectividades: Lucha a Volador (Defensores)}
\centering
\resizebox{\ifdim\width>\linewidth\linewidth\else\width\fi}{!}{
\fontsize{10}{12}\selectfont
\begin{tabular}[t]{>{}lrrrrrrrrr}
\toprule
\textbf{ } & \textbf{Lucha} & \textbf{Normal} & \textbf{Planta} & \textbf{Psíquico} & \textbf{Roca} & \textbf{Siniestro} & \textbf{Tierra} & \textbf{Veneno} & \textbf{Volador}\\
\midrule
\textbf{Acero} & 1.0 & 1 & 1.0 & 1.0 & 2.0 & 1.0 & 1.0 & 1.0 & 1.0\\
\textbf{Agua} & 1.0 & 1 & 0.5 & 1.0 & 2.0 & 1.0 & 2.0 & 1.0 & 1.0\\
\textbf{Bicho} & 0.5 & 1 & 2.0 & 2.0 & 1.0 & 2.0 & 1.0 & 0.5 & 0.5\\
\textbf{Dragón} & 1.0 & 1 & 1.0 & 1.0 & 1.0 & 1.0 & 1.0 & 1.0 & 1.0\\
\textbf{Eléctrico} & 1.0 & 1 & 0.5 & 1.0 & 1.0 & 1.0 & 0.0 & 1.0 & 2.0\\
\addlinespace
\textbf{Fantasma} & 1.0 & 0 & 1.0 & 2.0 & 1.0 & 0.5 & 1.0 & 1.0 & 1.0\\
\textbf{Fuego} & 1.0 & 1 & 2.0 & 1.0 & 0.5 & 1.0 & 1.0 & 1.0 & 1.0\\
\textbf{Hada} & 2.0 & 1 & 1.0 & 1.0 & 1.0 & 2.0 & 1.0 & 0.5 & 1.0\\
\textbf{Hielo} & 1.0 & 1 & 2.0 & 1.0 & 1.0 & 1.0 & 2.0 & 1.0 & 2.0\\
\textbf{Lucha} & 1.0 & 2 & 1.0 & 0.5 & 2.0 & 2.0 & 1.0 & 0.5 & 0.5\\
\addlinespace
\textbf{Normal} & 1.0 & 1 & 1.0 & 1.0 & 0.5 & 1.0 & 1.0 & 1.0 & 1.0\\
\textbf{Planta} & 1.0 & 1 & 0.5 & 1.0 & 2.0 & 1.0 & 2.0 & 0.5 & 0.5\\
\textbf{Psíquico} & 2.0 & 1 & 1.0 & 0.5 & 1.0 & 0.0 & 1.0 & 2.0 & 1.0\\
\textbf{Roca} & 0.5 & 1 & 1.0 & 1.0 & 1.0 & 1.0 & 0.5 & 1.0 & 2.0\\
\textbf{Siniestro} & 0.5 & 1 & 1.0 & 2.0 & 1.0 & 0.5 & 1.0 & 1.0 & 1.0\\
\addlinespace
\textbf{Tierra} & 1.0 & 1 & 0.5 & 1.0 & 2.0 & 1.0 & 1.0 & 2.0 & 0.0\\
\textbf{Veneno} & 1.0 & 1 & 2.0 & 1.0 & 0.5 & 1.0 & 0.5 & 0.5 & 1.0\\
\textbf{Volador} & 2.0 & 1 & 2.0 & 1.0 & 0.5 & 1.0 & 1.0 & 1.0 & 1.0\\
\bottomrule
\end{tabular}}
\end{table}

\chapter{Marco Metódico}

En esta sección se detallan los procedimientos algorítmicos, las
metodologías de procesamiento y las pruebas estadísticas. Se describe de
manera sistemática el tratamiento de la fuente de datos 29. Pokemon.csv,
aplicando técnicas de data wrangling (limpieza) y análisis descriptivo
para determinar la significancia de las diferencias entre las
estadísticas base de las tipologías elementales.

El análisis se ha desarrollado utilizando el lenguaje de programación R
bajo el entorno de desarrollo RStudio. Se ha adoptado el paradigma de
programación funcional ``Tidyverse'' para la manipulación eficiente de
grandes volúmenes de datos, asegurando un flujo de trabajo optimizado
que permite la generación de graficos comparativos de alta precisión
técnica.

\section{Bases metódicas utilizadas en el trabajo de investigación}

La presente investigación se define bajo un enfoque cuantitativo con un
alcance descriptivo-comparativo. Se centra en la recolección y el
análisis de datos numéricos preexistentes, permitiendo el uso de
herramientas estadísticas para identificar tendencias y patrones de
comportamiento de manera objetiva. Un alcance descriptivo, el estudio
busca caracterizar cada uno de los 18 tipos Pokémon según sus variables
de rendimiento. Finalmente, se integra un alcance comparativo que
permite contrastar las medias y varianzas entre los diversos tipos
elementales, estableciendo así las superioridades o debilidades
estadísticas que definen a cada categoría (Atacantes, Tanques, Muros).\\

\subsection{Procesamiento y Limpieza de Datos}

\begin{longtable}[]{@{}rrrrrrr@{}}
\caption{Resumen de datos faltantes por columna}\tabularnewline
\toprule\noalign{}
v\_tipo1 & v\_hp & v\_atk & v\_def & v\_sp\_atk & v\_sp\_def &
v\_speed \\
\midrule\noalign{}
\endfirsthead
\toprule\noalign{}
v\_tipo1 & v\_hp & v\_atk & v\_def & v\_sp\_atk & v\_sp\_def &
v\_speed \\
\midrule\noalign{}
\endhead
\bottomrule\noalign{}
\endlastfoot
0 & 0 & 0 & 0 & 0 & 0 & 0 \\
\end{longtable}

Descripción\\

\subsection{Item 2}

Descripción\\

\subsubsection{Item 2.2}

Descripción\\

\subsection{Item n}

Descripción\\

\chapter{Análisis de Resultados}

\{\small Este capítulo presenta los resultados obtenidos al aplicar las
metodologías descritas en el capítulo anterior, así como del análisis
estadístico descriptivo de las variables estudiadas y las ventajas y
desventajas al aplicar dichos tests.\}

\section{Presentación de resultados 1}

Descripción\\

\section{Presentación de resultados 2}

Descripción\\

\section{Presentación de resultados n}

Descripción\\

\chapter{Conclusiones y Recomendaciones}

\textbf{1.-} Conclusion 1 del trabajo de investigación.\\

\textbf{Recomendación} Recomendación 1 del trabajo de investigación.\\

\textbf{2.-} Conclusion 2 del trabajo de investigación.\\

\textbf{Recomendación} Recomendación 2 del trabajo de investigación.\\

\textbf{n.-} Conclusion n del trabajo de investigación.\\

\textbf{Recomendación} Recomendación n del trabajo de investigación.\\

\appendix
\renewcommand{\thechapter}{\Alph{chapter}}

\chapter{Anexo A}

\chapter{Anexo B}

\chapter{Anexo C}

%%%%%%%%%%%%%%%%%%%%%%%%%%%%%%%%%%%%%%%%%%%%%%%%%%%%%
%                   REFERENCIAS                     %
%%%%%%%%%%%%%%%%%%%%%%%%%%%%%%%%%%%%%%%%%%%%%%%%%%%%%
% Mantenemos la estructura natbib de tu plantilla
\renewcommand{\bibname}{Lista de Referencias}
\bibliographystyle{apalike} 
\bibliography{bibliografia.bib} 

%%%%%%%%%%%%%%%%%%%%%%%%%%%%%%%%%%%%%%%%%%%%%%%%%%%%%
%                   APÉNDICES                       %
%%%%%%%%%%%%%%%%%%%%%%%%%%%%%%%%%%%%%%%%%%%%%%%%%%%%%

\end{document}
